\begin{helper}
Let \(0<p<1\).  Fix a nonempty history event \(C\) in the two-bite sample
space, a representative sample \(\omega_0\in C\), a recorded coordinate set
\(R\), and a finite terminal candidate coordinate set \(U\).  Assume that
\(C\) is exactly the cylinder
\[
\omega\in C
\quad\Longleftrightarrow\quad
\omega\hbox{ has the same injection as }\omega_0
\hbox{ and }
\noderef{TwoBiteEdgeCoordinateValue}(\omega,e)
=\noderef{TwoBiteEdgeCoordinateValue}(\omega_0,e)
\hbox{ for every }e\in R.
\tag{1}
\]
Assume also that \(U\cap R=\emptyset\), and that every element of both \(R\)
and \(U\) is a red or blue non-loop base-edge coordinate represented with
increasing endpoint order.  This canonical-orientation hypothesis is part of
the independence assertion: it rules out the possibility that \(R\) records
\(\operatorname{red}(b,a)\) while \(U\) exposes
\(\operatorname{red}(a,b)\), or the analogous blue coordinate.  The set \(U\)
is only a fixed set of coordinates not constrained by the cylinder; it is not
assumed to be the final staged-open set.  Then, for every \(A\subseteq U\),
\[
\noderef{TwoBiteConditionalProbability}
\bigl(\text{exactly the coordinates in }A\text{ are present on }U
\mid C\bigr)
=p^{|A|}(1-p)^{|U|-|A|}.
\tag{2}
\]
\end{helper}
\begin{proof}
Let \(\mathcal U\) be the finite tagged coordinate universe consisting of all
red and blue increasing non-loop base-edge coordinates in \(\mathrm{Fin}(m)\).
The hypotheses put \(R\subseteq\mathcal U\) and \(U\subseteq\mathcal U\).
Write
\[
W=\mathcal U\setminus(R\cup U).
\]
Since \(R\) and \(U\) are both subsets of \(\mathcal U\) and \(U\cap R=\emptyset\),
the three finite sets \(R\), \(U\), and \(W\) are disjoint and cover every red
or blue edge coordinate whose
\noderef{GnpGraphWeight} factor can vary after the injection is fixed.
Here the phrase ``can vary'' is literal after unfolding the weight
definitions.  A two-bite sample is a red graph, a blue graph, and an
injection.  By \(\noderef{TwoBiteSampleWeight}\), its weight is
\[
\noderef{GnpGraphWeight}(p,G_R)\,
\noderef{GnpGraphWeight}(p,G_B)\,
\noderef{UniformInjectionWeight}(n,m).
\]
The injection factor is independent of every red and blue edge value.  Each
\(\noderef{GnpGraphWeight}\) is the finite product over increasing non-loop
base pairs, with factor \(p\) when the corresponding edge is present and
factor \(1-p\) when it is absent.  Finally,
\(\noderef{TwoBiteEdgeCoordinateValue}\) simply reads the red adjacency for a
red tagged coordinate and the blue adjacency for a blue tagged coordinate.
Thus the graph-coordinate factors relevant to this proof are exactly the
tagged increasing non-loop coordinates in \(\mathcal U\).

The cylinder description (1) has two consequences.  First, every sample in
\(C\) has the same injection as \(\omega_0\).  Second, every coordinate in
\(R\) has the same truth value as in \(\omega_0\).  Conversely, there are no
other restrictions in the definition of \(C\): if the injection and the
\(R\)-table agree with \(\omega_0\), then the sample is in \(C\).  In
particular, the coordinates in \(U\) and \(W\) are unrestricted by the
conditioning event.

The last sentence uses the canonical-orientation hypotheses in an essential
way.  Suppose \(e\in U\) and a recorded coordinate \(r\in R\) reads the same
underlying graph adjacency as \(e\).  The two coordinates must have the same
color, because red and blue graph factors are distinct.  If
\(e=\operatorname{red}(a,b)\) and \(r=\operatorname{red}(c,d)\), then both
\(a<b\) and \(c<d\).  Equality of the underlying unordered base edge therefore
forces \(a=c\) and \(b=d\), so \(r=e\); the blue case is identical.  This
contradicts \(U\cap R=\emptyset\).  Thus no recorded coordinate, even through
the symmetry of simple-graph adjacency, fixes a coordinate of \(U\).

Because \(\omega_0\in C\), the cylinder has at least one sample.  The
injection factor
\(\noderef{UniformInjectionWeight}(n,m)\) is positive on any injection that
occurs in the sample space, and \(0<p<1\) makes every factor \(p\) and
\(1-p\) in the red and blue graph weights positive.  More explicitly, the
injection of \(\omega_0\) exists, so the number of injections is nonzero and
the reciprocal defining
\(\noderef{UniformInjectionWeight}(n,m)\) is positive.  The fixed
\(R\)-truth table contributes a finite product of factors, each equal to
either \(p\) or \(1-p\), hence a positive factor.  The coordinates in \(U\)
and \(W\) are not restricted by the cylinder, so they contribute positive
total mass.  Therefore
\[
\noderef{TwoBiteEventProbability}(C)>0,
\tag{3}
\]
so the nonzero branch of
\(\noderef{TwoBiteConditionalProbability}\) is used.

Fix \(A\subseteq U\).  In the numerator for the event that exactly \(A\)
occurs on \(U\), unfold \(\noderef{TwoBiteEventProbability}\).  It is the
finite sum of \(\noderef{TwoBiteSampleWeight}\) over those samples which both
lie in \(C\) and satisfy the \(U\)-event.  For every summand, the injection
is the fixed injection of \(\omega_0\), the \(R\)-coordinates have the fixed
table inherited from \(\omega_0\), and the \(U\)-coordinates are fixed by
\[
e\in A \Longleftrightarrow
\noderef{TwoBiteEdgeCoordinateValue}(\omega,e).
\]
Because \(A\subseteq U\), every coordinate of \(U\) contributes \(p\) exactly
when it lies in \(A\), and contributes \(1-p\) exactly when it lies in
\(U\setminus A\).  Thus the contribution of \(U\) to every numerator summand is
\[
p^{|A|}(1-p)^{|U|-|A|}. \tag{4}
\]
The coordinates in \(W\) are not mentioned either by the cylinder \(C\) or by
the event on \(U\).  Summing over the two possible values of any one
\(w\in W\) gives \(p+(1-p)=1\).  Iterating over the finite set \(W\), all
\(W\)-coordinates contribute a total factor \(1\).  In Lean this last
iteration is the usual finite-product calculation using
\(\texttt{Finset.prod\_eq\_one}\) after each unrestricted coordinate has
been summed out.

The denominator \(\noderef{TwoBiteEventProbability}(C)\) has the same fixed
injection factor and the same fixed \(R\)-table factor as the numerator,
because the denominator sums over the same cylinder \(C\) but imposes no
condition on \(U\).  Its \(W\)-coordinates again sum to \(1\).  Its
\(U\)-coordinates are unrestricted, so their total
factor is
\[
\sum_{B\subseteq U}p^{|B|}(1-p)^{|U|-|B|}
=\prod_{e\in U}(p+(1-p))=1. \tag{5}
\]
Consequently the numerator is exactly the common fixed injection and
recorded-table factor multiplied by the \(U\)-factor in (4), while the
denominator is the same common factor multiplied by \(1\).  The common factor
is positive by (3), and
\(\noderef{TwoBiteConditionalProbability}\) is therefore the ordinary quotient
of numerator by denominator.  Cancelling the common factor and dividing by
the positive denominator gives (2).
\end{proof}
